\documentclass[11pt]{article}
\usepackage[margin=2.1cm]{geometry}
\usepackage{booktabs}
\usepackage{amsmath}
\usepackage{array}
\usepackage{xcolor}

\newcommand{\blue}[1]{\textcolor{blue}{#1}}
\newcommand{\cobs}{c_{\mathrm{obs}}}
\newcommand{\cbal}{c_{\mathrm{bal}}}
\newcommand{\eroot}{\varepsilon_{\mathrm{root}}}

\title{KRegCCD: the 15 four-taxon tree probabilities, symbolically}
\author{}
\date{}

\begin{document}
\maketitle
\vspace{-2.5em}

\section*{Setup}

Taxa $\{A,B,C,D\}$. Two trees observed, each once:
\[
T_1 = (((A,B),C),D), \qquad T_2 = (((D,C),B),A),
\]
so the observed non-trivial clades are $\{A,B\},\{A,B,C\}$ (from $T_1$) and
$\{C,D\},\{B,C,D\}$ (from $T_2$). Clades never seen are the \blue{blue nodes}.
KRegCCD is fitted with smoothing parameter $\alpha$ and per-clade escape
probability $\mu$ (reserve depth $k=2$, \texttt{NovelMode.FLAT}). With four taxa
there are exactly $(2\cdot4-3)!!=15$ rooted topologies, so this is the whole
tree space.

\section*{The three building blocks}

The escape probability $\mu$ is a single global parameter: every reservable
clade $C$ reserves the \emph{same} total mass $\mu$ for its blue resolutions. The
escape root $\varepsilon(C)$, by contrast, is solved \emph{per clade} from that
clade's own reserve equation
\[
R(C) \;=\; \sum_{j\ge 1} N_j(C)\,\varepsilon(C)^{\,j} \;=\; \mu ,
\]
where $N_j(C)$ counts the depth-$j$ novel resolutions of $C$. So $\varepsilon(C)$
varies from clade to clade (the root and the 3-taxon clades below get different
roots), while $\mu$ does not. The two distinct values $\eroot$ and
$\varepsilon_3$ used in the table are $\varepsilon(C)$ at the root and at a
3-taxon clade respectively.

\paragraph{Backbone CCPs (from $\alpha$).}
Additive-$\alpha$ smoothing at the root clade $\{A,B,C,D\}$, which was seen twice
and carries three candidate splits (the two observed splits plus the
split-expanded balanced split $\{A,B\}|\{C,D\}$), gives
\[
\cobs \;=\; \frac{1+\alpha}{2+3\alpha}
\quad(\text{each observed split}),
\qquad
\cbal \;=\; \frac{\alpha}{2+3\alpha}
\quad(\text{balanced split}),
\]
with the key identity $\;2\cobs+\cbal = 1$. The two interior 3-taxon clades each
carry a single split, so their CCP is $1$.

\paragraph{Root escape root $\eroot$.}
The root reserves total mass $\mu$ for its blue resolutions. In FLAT mode the
reserve sums tree counts: $M_1=2$ trees carry one novel clade and $M_2=6$ carry
two, so the root's escape root $\eroot$ solves
\[
2\eroot + 6\eroot^2 = \mu
\qquad\Longrightarrow\qquad
\eroot \;=\; \frac{\sqrt{1+6\mu}-1}{6}.
\]

\paragraph{Interior 3-clade escape root $\varepsilon_3\equiv\varepsilon(\text{3-clade})$.}
A 3-taxon observed clade has two blue resolutions and no deeper order, so its
escape root solves $2\varepsilon_3=\mu$, i.e. $\varepsilon_3=\mu/2$. Same $\mu$ as
the root, but a different root because the $N_j$ counts differ.

\section*{The 15 probabilities}

Each tree is a product over its clades: a reservable observed clade that keeps an
observed split contributes $(1-\mu)\,\mathrm{CCP}$; a clade that escapes
contributes $\varepsilon(C)^{\,\#\text{novel}}$ for its blue region; cherries
contribute $1$.

\begin{center}
\renewcommand{\arraystretch}{1.35}
\begin{tabular}{c l l c r}
\toprule
 & \textbf{Newick} & \textbf{Category} & \textbf{Symbolic} & \textbf{Value} \\
\midrule
$T_1$ & $(((A,B),C),D)$ & observed            & $(1-\mu)^2\,\cobs$ & $0.39484375$ \\
$T_2$ & $(((D,C),B),A)$ & observed            & $(1-\mu)^2\,\cobs$ & $0.39484375$ \\
$T_3$ & $((A,B),(C,D))$ & $\alpha$-expansion  & $(1-\mu)\,\cbal$   & $0.11875000$ \\
\midrule
\multicolumn{5}{l}{\emph{One blue node: a 3-taxon clade over an observed cherry $(\#\text{novel}=1$, root escapes$)$}}\\
4 & $(((A,B),D),C)$ & blue \blue{$\{A,B,D\}$} & $\eroot$ & $0.02336257$ \\
5 & $(((C,D),A),B)$ & blue \blue{$\{A,C,D\}$} & $\eroot$ & $0.02336257$ \\
\midrule
\multicolumn{5}{l}{\emph{One blue node: a novel cherry under an observed triple $($root stays observed, 3-clade escapes$)$}}\\
6 & $(((A,C),B),D)$ & blue \blue{$\{A,C\}$} & $(1-\mu)\,\cobs\,\dfrac{\mu}{2}$ & $0.01039063$ \\
7 & $(((B,C),A),D)$ & blue \blue{$\{B,C\}$} & $(1-\mu)\,\cobs\,\dfrac{\mu}{2}$ & $0.01039063$ \\
8 & $(((B,C),D),A)$ & blue \blue{$\{B,C\}$} & $(1-\mu)\,\cobs\,\dfrac{\mu}{2}$ & $0.01039063$ \\
9 & $(((B,D),C),A)$ & blue \blue{$\{B,D\}$} & $(1-\mu)\,\cobs\,\dfrac{\mu}{2}$ & $0.01039063$ \\
\midrule
\multicolumn{5}{l}{\emph{Two blue nodes: fully novel $(\#\text{novel}=2$, root escapes$)$}}\\
10 & $(((A,D),B),C)$ & blue \blue{$\{A,D\},\{A,B,D\}$} & $\eroot^2$ & $0.00054581$ \\
11 & $(((B,D),A),C)$ & blue \blue{$\{B,D\},\{A,B,D\}$} & $\eroot^2$ & $0.00054581$ \\
12 & $(((A,C),D),B)$ & blue \blue{$\{A,C\},\{A,C,D\}$} & $\eroot^2$ & $0.00054581$ \\
13 & $(((A,D),C),B)$ & blue \blue{$\{A,D\},\{A,C,D\}$} & $\eroot^2$ & $0.00054581$ \\
14 & $((A,C),(B,D))$ & blue \blue{$\{A,C\},\{B,D\}$}   & $\eroot^2$ & $0.00054581$ \\
15 & $((A,D),(B,C))$ & blue \blue{$\{A,D\},\{B,C\}$}   & $\eroot^2$ & $0.00054581$ \\
\midrule
   &                 & \textbf{sum}                   & $1$ & $1.00000000$ \\
\bottomrule
\end{tabular}
\end{center}

\noindent Values shown at the fitted $\alpha=0.40,\ \mu=0.05$, for which
$\cobs=\tfrac{7}{16}=0.4375$, $\cbal=\tfrac18=0.125$, and
$\eroot=\tfrac{\sqrt{1.3}-1}{6}=0.0233625708\ldots$

\section*{Why it sums to exactly 1}

Group the fifteen terms by the three building blocks:
\begin{align*}
\sum_{\text{all }15}
&= \underbrace{2(1-\mu)^2\cobs}_{T_1,T_2}
 + \underbrace{(1-\mu)\cbal}_{T_3}
 + \underbrace{4(1-\mu)\cobs\tfrac{\mu}{2}}_{\text{trees }6\text{--}9}
 + \underbrace{2\eroot + 6\eroot^2}_{\text{trees }4,5,10\text{--}15} \\[2pt]
&= (1-\mu)\,\cobs\big[\,2(1-\mu)+2\mu\,\big] + (1-\mu)\cbal + \big(2\eroot+6\eroot^2\big)\\[2pt]
&= (1-\mu)\big(2\cobs+\cbal\big) + \mu
 \;=\; (1-\mu)\cdot 1 + \mu \;=\; 1.
\end{align*}
The escaped mass $2\eroot+6\eroot^2$ is, by definition of $\eroot$,
exactly the reserve $\mu$; and the backbone CCPs satisfy $2\cobs+\cbal=1$. So the
distribution is normalised for any $\alpha\in(0,\infty)$ and $\mu\in(0,1)$, with
no leftover reserve. (This holds under \texttt{TailMode.NONE}. The optimiser's
default \texttt{TailMode.BOUND} subtracts a strict upper bound on the
deeper-than-$k$ reserve tail, which is empty here, so it sub-normalises to
$0.99975$; on a finite enumerated space \texttt{NONE} is exact.)

\end{document}
